\subsection{The monodromy condition and the Fontaine comparison}

\begin{proposition}[B-RD4: Lead B' needs an additional hypothesis]\label{prop:B-RD4}
\status{Refuted as stated.}
For every odd prime $p$, there is a height-two Barsotti--Tate group over
$K=\mathbf Q_p(\zeta_p)$ satisfying the manuscript's monodromy condition
whose rational Tate module is not potentially Hodge--Tate, and therefore
is not potentially semistable. Its manuscript inertia character is the
constant function $2$. Thus the monodromy condition in 9:11 cannot be
identified with potential semistability, and Fontaine's usual
Weil--Deligne representation need not be defined for its Tate module.
\end{proposition}
\begin{proof}
Write $\chi:G_K\to\mathbf Z_p^\times$ for the cyclotomic character.
Since $\mu_p\subset K$, it takes its values in $1+p\mathbf Z_p$. Put
\[
 \rho(g)=\begin{pmatrix}1&p^{-1}\log\chi(g)\\0&1\end{pmatrix}.
\]
This is a continuous representation on $T=\mathbf Z_p^2$, because the
logarithm is additive on $1+p\mathbf Z_p$ and takes values in
$p\mathbf Z_p$. Over a characteristic-zero field, the equivalence
between finite etale group schemes and finite continuous Galois modules
applied to $T/p^nT$ gives compatible finite etale group schemes $M[p^n]$.
The transition maps are induced by multiplication by $p$ on the lattices,
and the resulting inductive system is a height-two Barsotti--Tate group.
The invariant first coordinate and the quotient both give the constant
group $\mathbf Q_p/\mathbf Z_p$. They extend over $\mathcal O_K$, so the
specified two-step filtration satisfies the manuscript condition and
its character is $1+1=2$.

For $V=T\otimes\mathbf Q_p$, the Sen operator is
\[
 \Theta_V=\frac1p\begin{pmatrix}0&1\\0&0\end{pmatrix}.
\]
Indeed the logarithm of a unipotent matrix $1+aN$, with $N^2=0$, is
$aN$, and division by $\log\chi(g)$ gives the displayed operator.
It remains nonzero and nilpotent after restriction to any finite
extension. Sen's criterion says that a representation is Hodge--Tate
exactly when this operator is diagonalizable with integral eigenvalues;
the identical unscaled example and the criterion are explicitly given
in \cite[\S II.1.2, p.~9 of the author's version]{Berger2004Introduction}.
The operator is not diagonalizable. Potentially semistable
representations are potentially Hodge--Tate, so no finite extension
makes this representation semistable.
\end{proof}

\begin{remark}[The corrected scope of the proposed dictionary]
If one separately assumes potential semistability of $T_p(M)\otimes
\mathbf Q_p$, the comparison proposed in Lead B' is the appropriate
one: on a potentially good graded piece, crystalline comparison
identifies its Dieudonne isocrystal with the corresponding period
module, equivariantly for descent. Exactness of the potentially
semistable period functor then adds these characters across a
filtration. Covariant and contravariant conventions must be fixed:
with $D_{\mathrm{cov}}(G)=D_{\mathrm{contra}}(G^D)$, the covariant
isocrystal corresponds to $D_{\mathrm{cris}}(T_pG)(-1)$; the Tate twist
has trivial Weil--Deligne inertia character. This paragraph records
the corrected comparison to be sourced; the counterexample above
requires no comparison theorem.
\end{remark}

\begin{proposition}[B-RD5: the height-one case of B.4]\label{prop:B-RD5}
\status{Proved.}
Suppose $p\geq3$ and $\mu_p\subset K$. For height $h=1$, the manuscript
character on prime-to-$p$ inertia elements is determined by $M(0)$.
The least index in its indexing is therefore $\nu=0$.
\end{proposition}
\begin{proof}
A height-one group satisfying the monodromy condition itself acquires
good reduction. Its good model has dimension zero or one. In dimension
zero it is etale, and in dimension one its Cartier dual is etale.
Over the completed maximal unramified extension of a good-reduction
field its model is respectively $\mathbf Q_p/\mathbf Z_p$ or
$\mu_{p^\infty}$. Its finite descent inertia character therefore takes
values in the torsion subgroup of $\mathbf Z_p^\times$, namely
$\mu_{p-1}$. In dimension zero, reduction modulo $p$ gives precisely
this character. In dimension one, the Tate module acquires additionally
the cyclotomic factor, whose reduction is trivial under $\mu_p\subset K$.
Thus in both cases the finite character is the unique Teichmuller lift
of the one-dimensional mod-$p$ inertia character. The same conclusion
holds after summing characters of height-one potentially good graded
pieces: the trace on a prime-to-$p$ element is the Brauer character of
the mod-$p$ representation, and is independent of the chosen filtration.
\end{proof}
